%!TEX root = ../main.tex
% Final LOGIKA UI 2026 - Geometri
\begin{center}
    {\color{crimson}\textbf{\Huge Final LOGIKA UI 2026: Geometri}}\\[4pt]
    {\color{inksoft}\textbf{\large Problem and Solution proposed by Azzam L. H.}}
\end{center}

\section{Soal}
Diberikan segitiga lancip $ABC$ dengan $\angle BAC > \angle ACB$ dan pusat lingkaran luar $O$. Misalkan $O_1$ dan $O_2$ berturut-turut adalah pusat lingkaran luar segitiga $ABO$ dan $ACO$. Lingkaran luar segitiga $AO_1O_2$ memotong sisi $BC$ di dua titik berbeda $P$ dan $Q$ dengan urutan $B, P, Q, C$. Misalkan $O_3$ adalah pusat lingkaran luar segitiga $OPQ$, dan garis sumbu ruas garis $OB$ memotong garis $O_2O_3$ di titik $D$. Buktikan bahwa
\[
    AO_3 = O_3D \quad\text{jika dan hanya jika}\quad A,\ O_1,\ D \text{ segaris.}
\]

% Post 51 - figure for the problem statement.
% Exact configuration A = 76, B = 74, C = 30 (degrees).
%!TEX root = ../main.tex
% Shared colors for the post 51 figures (Helvethia palette + teal/orange accents).
\colorlet{figmain}{ink}
\colorlet{figaux}{inksoft}
\colorlet{figcircle}{inksoft!45}
\definecolor{figteal}{RGB}{25,130,115}
\colorlet{figorange}{orange!80!black}

\begin{center}
\begin{tikzpicture}[scale=0.45, line cap=round, line join=round]
    \clip (-7.6, -12.4) rectangle (9.2, 11.1);

    \coordinate (O)  at (  0.0000,   0.0000);
    \coordinate (A)  at ( -4.3069,   4.4599);
    \coordinate (B)  at ( -6.0158,  -1.4999);
    \coordinate (C)  at (  6.0158,  -1.4999);
    \coordinate (O1) at ( -3.4409,   0.9867);
    \coordinate (O2) at (  5.6233,   9.7399);
    \coordinate (O3) at (  2.1824,  -2.2600);
    \coordinate (P)  at ( -0.8660,  -1.4999);
    \coordinate (Q)  at (  5.2308,  -1.4999);
    \coordinate (D)  at ( -0.3925, -11.2398);
    \coordinate (K)  at (  2.1824,   4.2334);

    % (ABC), (AO1O2), (OPQ)
    \draw[color=figcircle, line width=0.55pt] (O) circle (6.2000);
    \draw[color=figorange, line width=0.55pt, dashed] (K) circle (6.4932);
    \draw[color=figteal, line width=0.80pt] (O3) circle (3.1417);

    % garis sumbu OB (melalui O1 dan D) dan garis O2O3
    \draw[color=figaux, line width=0.6pt] ($(D)!-0.03!(O1)$) -- ($(D)!1.30!(O1)$);
    \draw[color=figaux, line width=0.6pt] ($(O2)!-0.03!(D)$) -- ($(O2)!1.03!(D)$);
    \draw[color=figaux, line width=0.5pt] (O) -- (B);

    \draw[color=figmain, line width=1.1pt] (A) -- (B) -- (C) -- cycle;

    \fill[figmain] (A)  circle (2.2pt) node[above left]  {$A$};
    \fill[figmain] (B)  circle (2.2pt) node[below left]  {$B$};
    \fill[figmain] (C)  circle (2.2pt) node[below right] {$C$};
    \fill[figmain] (O)  circle (2.2pt) node[above right] {$O$};
    \fill[crimson] (O1) circle (2.2pt) node[left, xshift=-2pt] {$O_1$};
    \fill[crimson] (O2) circle (2.2pt) node[above left]  {$O_2$};
    \fill[crimson] (O3) circle (2.2pt) node[right, xshift=2pt, yshift=-2pt] {$O_3$};
    \fill[figmain] (P)  circle (2.0pt) node[below left]  {$P$};
    \fill[figmain] (Q)  circle (2.0pt) node[below, xshift=-4pt] {$Q$};
    \fill[crimson] (D)  circle (2.2pt) node[below left]  {$D$};
\end{tikzpicture}
\end{center}


\newpage
\section{Motivasi dan Cerita Terbentuknya Soal Ini}
Kalo ngga salah sekitar bulan Agustus-September 2026 ini, aku lagi iseng \textit{scrolling} AoPS, soalnya aku juga lagi ngescrapping bank-bank soal buat dimasukin ke \url{helvethia.web.id}. Nah, ternyata di AoPS, soal-soal Jepang sempat di"hide" oleh moderator (karena \textit{copyright problem} (?)).

Yaudah, aku cari lah di sumber aslinya \url{imojp.org} dan akhirnya nemu. Singkat cerita, setelah ngescrap dan ngumpulin jadi bank soal, tentu saja aku coba-coba lah dari soal yang tahunnya agak baru. Karena aku lihat yang tahun 2026 dan ngga terlalu suka, akhirnya setelah nyoba soal geometrinya \textcrimson{\textbf{Japan MO 2025 Final}}, aku langsung suka. 

Ya, random aja pas nemu ini kayak, "bro, ini soal apaan aneh juga" : $A$, $O$, $O_3$ segaris.

Yaudahlah, karena emang kebetulan sekitar September-Oktober diminta jadi tim soal LOGIKA UI 2026, aku coba lah ulik ini soal. Singkat cerita (time skip ke awal November), setelah agak berdebu, aku nemu lagi bank soal yang aku tulis di buku tulis, dan sialnya lupa aku tulis sumber tahunnya (\textit{but I know it's an Iranian MO problem, maybe?}) yang mana ngga aku catat tahun berapa, terlalu malas aku ceknya itu tahun berapa. Dari situ, aku coba lah idenya (perpotongan garis dan lingkaran) dan seperti peyuka soal geometri pada umumnya, aku teruskan dengan main-main di GeoGebra. Dan ya, setelah nemu momen "aha"nya jadilah, soal cantik tapi mematikan ini. (No, Andrew adalah anomali, mana ada orang normal bisa solve ini 5 menit broo x\_x).

\newpage
\section{Lemma Inspirasi Utama Soal Ini (Japan MO 2025 Final)}
\begin{lemma*}
    Diberikan segitiga lancip $ABC$ dengan pusat lingkaran luar $O$. Misalkan $O_1$ dan $O_2$ berturut-turut adalah pusat lingkaran luar segitiga $ABO$ dan $ACO$. Lingkaran luar segitiga $AO_1O_2$ memotong sisi $BC$ di dua titik berbeda $P$ dan $Q$ dengan urutan $B, P, Q, C$. Jika $O_3$ adalah pusat lingkaran luar segitiga $OPQ$, maka titik $A$, $O$, dan $O_3$ segaris. Lebih jauh, titik $P$ terletak pada lingkaran luar segitiga $ABO$, titik $Q$ terletak pada lingkaran luar segitiga $ACO$, dan $O$ adalah pusat lingkaran dalam segitiga $APQ$.
\end{lemma*}

% Post 51 - figure for the JMO 2025 Final lemma.
%!TEX root = ../main.tex
% Shared colors for the post 51 figures (Helvethia palette + teal/orange accents).
\colorlet{figmain}{ink}
\colorlet{figaux}{inksoft}
\colorlet{figcircle}{inksoft!45}
\definecolor{figteal}{RGB}{25,130,115}
\colorlet{figorange}{orange!80!black}

\begin{center}
\begin{tikzpicture}[scale=0.50, line cap=round, line join=round]
    \coordinate (A)  at (-7.572, 11.157);
    \coordinate (B)  at (-14.255, 0.689);
    \coordinate (C)  at (1.678, 0.818);
    \coordinate (O)  at (-6.307, 2.982);
    \coordinate (O1) at (-11.584, 6.350);
    \coordinate (O2) at (-0.650, 8.042);
    \coordinate (O3) at (-5.834, -0.077);
    \coordinate (P)  at (-8.821, 0.733);
    \coordinate (Q)  at (-2.860, 0.781);

    \draw[color=figcircle, line width=0.5pt] (-5.880, 5.667) circle (5.744);
    \draw[color=figorange, line width=0.5pt, dashed] (O1) circle (6.261);
    \draw[color=figorange, line width=0.5pt, dashed] (O2) circle (7.590);
    \draw[color=figteal,   line width=0.75pt] (O3) circle (3.095);

    \draw[color=figaux, line width=0.55pt] (A) -- (O);
    \draw[color=figaux, line width=0.55pt] (O) -- (C);
    \draw[color=figaux, line width=0.55pt] (B) -- (O);
    \draw[color=figaux, line width=0.55pt] (P) -- (A);
    \draw[color=figaux, line width=0.55pt] (A) -- (Q);
    \draw[color=figaux, line width=0.55pt] (O1) -- (Q);
    \draw[color=figaux, line width=0.55pt] (P) -- (O2);
    \draw[color=figaux, line width=0.55pt] (O1) -- (O2);
    \draw[color=figaux, line width=0.55pt] (O3) -- (P);
    \draw[color=crimson, line width=0.85pt] ($(A)!-0.06!(O3)$) -- ($(A)!1.10!(O3)$);

    \draw[color=figmain, line width=1.1pt] (A) -- (B) -- (C) -- cycle;

    \fill[figmain] (A)  circle (2.2pt) node[above right] {$A$};
    \fill[figmain] (B)  circle (2.2pt) node[below left]  {$B$};
    \fill[figmain] (C)  circle (2.2pt) node[below right] {$C$};
    \fill[crimson] (O)  circle (2.4pt) node[above right] {$O$};
    \fill[figmain] (O1) circle (2.2pt) node[above left]  {$O_1$};
    \fill[figmain] (O2) circle (2.2pt) node[above right] {$O_2$};
    \fill[crimson] (O3) circle (2.4pt) node[below right] {$O_3$};
    \fill[figmain] (P)  circle (2.0pt) node[below, yshift=-2pt] {$P$};
    \fill[figmain] (Q)  circle (2.0pt) node[below, yshift=-2pt] {$Q$};
\end{tikzpicture}
\end{center}


\newpage
\begin{proof}[\textbf{Klaim 1.}] Titik $P$ terletak pada lingkaran luar $(ABO)$ dan titik $Q$ terletak pada lingkaran luar $(ACO)$, serta $O$ merupakan pusat lingkaran dalam $\triangle APQ$.\\
    \textbf{Bukti Klaim.} Karena $O_1$ adalah pusat $(ABO)$ dan $O_2$ adalah pusat $(ACO)$, garis hubung pusat $O_1O_2$ merupakan garis sumbu dari $OA$ ($O_1O_2 \perp OA$). Misalkan $P'$ adalah titik potong kedua garis $BC$ dengan $(ABO)$. Menggunakan sudut berarah modulo $180^\circ$:
    \begin{align*}
        \measuredangle(O_1A, AP') &= 90^\circ + \measuredangle(BA, BC), \\
        \measuredangle(OA, OP') &= \measuredangle(BA, BP') = \measuredangle(BA, BC) = 90^\circ - \measuredangle(CO, CA) = \measuredangle(OA, OO_2),
    \end{align*}
    sehingga $O, P', O_2$ segaris. Akibatnya,
    \[
        \measuredangle(O_1O_2, O_2P') = \measuredangle(O_1O_2, OA) + \measuredangle(OA, OO_2) = 90^\circ + \measuredangle(BA, BC) = \measuredangle(O_1A, AP'),
    \]
    yang berarti $P' \in (AO_1O_2)$. Secara simetris, titik potong kedua $Q'$ dari $BC$ dengan $(ACO)$ juga terletak pada $(AO_1O_2)$, sehingga $P = P' \in (ABO)$ dan $Q = Q' \in (ACO)$.

    Selanjutnya, karena $OB = OC$ dan $B, P, Q, C$ segaris, berlaku
    \[
        \measuredangle PAO = \measuredangle PBO = \measuredangle OCQ = \measuredangle OAQ \quad\text{dan}\quad \measuredangle OPA = \measuredangle OBA = \measuredangle BAO = \measuredangle OPQ.
    \]
    Dengan demikian, $O$ terletak pada garis bagi dalam sudut $\angle PAQ$ maupun $\angle APQ$, sehingga $O$ adalah pusat lingkaran dalam $\triangle APQ$. Klaim 1 terbukti.
\end{proof}
\hrulefill
\begin{proof}[\textbf{Klaim 2.}] Titik-titik $A$, $O$, dan $O_3$ segaris.\\
    \textbf{Bukti Klaim.} Karena $O$ adalah pusat lingkaran dalam $\triangle APQ$ dan $O_3$ adalah pusat lingkaran luar $\triangle OPQ$, berdasarkan Lemma \textit{Incenter-Excenter} pada $\triangle APQ$, titik $O_3$ merupakan titik tengah busur $PQ$ (yang tidak memuat $A$) pada lingkaran luar $\triangle APQ$. Akibatnya, garis bagi dalam $AO$ dari $\angle PAQ$ melalui titik $O_3$, sehingga $A$, $O$, dan $O_3$ segaris. Klaim 2 terbukti.
\end{proof}
Dapat disimpulkan dari Klaim 2, lemma "Lemma Jepang" ini terbukti.

\newpage
\section{Solusi}
Dari \textcrimson{\textbf{Lemma Jepang tersebut}}, $A$, $O$, $O_3$ segaris dengan $O$ di antara $A$ dan $O_3$. Karena $O_1O = O_1B$ dan $DO = DB$, garis $O_1D$ adalah garis sumbu $OB$. Karena $O_2O = O_2Q$ (dari segiempat siklis $AOQC$) dan $O_3O = O_3Q$, garis $O_2O_3$ (yakni $O_3D$) adalah garis sumbu $OQ$.

% Post 51 - figure for the solution.
% Exact configuration A = 76, B = 74, C = 30 (degrees), where
% A, O1, D are collinear and AO3 = O3D.
%!TEX root = ../main.tex
% Shared colors for the post 51 figures (Helvethia palette + teal/orange accents).
\colorlet{figmain}{ink}
\colorlet{figaux}{inksoft}
\colorlet{figcircle}{inksoft!45}
\definecolor{figteal}{RGB}{25,130,115}
\colorlet{figorange}{orange!80!black}

\begin{center}
\begin{tikzpicture}[scale=0.40, line cap=round, line join=round]
    \clip (-7.5, -12.3) rectangle (7.5, 10.6);

    \coordinate (O)  at (  0.0000,   0.0000);
    \coordinate (A)  at ( -4.3069,   4.4599);
    \coordinate (B)  at ( -6.0158,  -1.4999);
    \coordinate (C)  at (  6.0158,  -1.4999);
    \coordinate (O1) at ( -3.4409,   0.9867);
    \coordinate (O2) at (  5.6233,   9.7399);
    \coordinate (O3) at (  2.1824,  -2.2600);
    \coordinate (P)  at ( -0.8660,  -1.4999);
    \coordinate (Q)  at (  5.2308,  -1.4999);
    \coordinate (D)  at ( -0.3925, -11.2398);

    \draw[color=figcircle, line width=0.55pt] (O)  circle (6.2000);
    \draw[color=figteal,   line width=0.80pt] (O3) circle (3.1417);
    \draw[color=figorange, line width=0.5pt, dashed] (O1) circle (3.5796);
    \draw[color=figorange, line width=0.5pt, dashed] (O2) circle (11.2467);

    \draw[color=figaux, line width=0.5pt] (O) -- (A);
    \draw[color=figaux, line width=0.5pt] (O) -- (B);
    \draw[color=figaux, line width=0.5pt] (O) -- (C);
    \draw[color=figaux, line width=0.5pt] (O) -- (P);
    \draw[color=figaux, line width=0.5pt] (O) -- (Q);
    \draw[color=figaux, line width=0.5pt] (O1) -- (O);
    \draw[color=figaux, line width=0.5pt] (O1) -- (O3);

    % A - O - O3 segaris
    \draw[color=crimson, line width=0.85pt] ($(A)!-0.08!(O3)$) -- ($(A)!1.12!(O3)$);
    % A - O1 - D (garis sumbu OB)
    \draw[color=crimson, line width=0.85pt] ($(A)!-0.06!(D)$) -- ($(A)!1.05!(D)$);
    % garis sumbu OQ (melalui O2, O3, D)
    \draw[color=figteal, line width=0.85pt] ($(O2)!-0.03!(D)$) -- ($(O2)!1.04!(D)$);

    \draw[color=figmain, line width=1.1pt] (A) -- (B) -- (C) -- cycle;

    \fill[figmain] (A)  circle (2.2pt) node[above left]  {$A$};
    \fill[figmain] (B)  circle (2.2pt) node[below left]  {$B$};
    \fill[figmain] (C)  circle (2.2pt) node[below right] {$C$};
    \fill[figmain] (O)  circle (2.2pt) node[above right] {$O$};
    \fill[crimson] (O1) circle (2.2pt) node[left, xshift=-2pt] {$O_1$};
    \fill[figmain] (O2) circle (2.0pt) node[above left]  {$O_2$};
    \fill[crimson] (O3) circle (2.2pt) node[right, xshift=2pt, yshift=-1pt] {$O_3$};
    \fill[figmain] (P)  circle (2.0pt) node[below left]  {$P$};
    \fill[figmain] (Q)  circle (2.0pt) node[below, xshift=-4pt] {$Q$};
    \fill[crimson] (D)  circle (2.2pt) node[below left]  {$D$};
\end{tikzpicture}
\end{center}


\begin{proof}[\textbf{Klaim 1.}] Segiempat $OO_1O_3D$ siklis dan $\angle O_1AO_3 = \angle O_1DO_3 = C$.\\
    \textbf{Bukti Klaim.} Dari segiempat siklis $AOQC$ diperoleh $\angle OQB = 180^\circ - \angle OQC = \angle OAC = 90^\circ - B$, sehingga dengan $\angle OBQ = 90^\circ - A$, jumlah sudut $\triangle OBQ$ menyebabkan $\angle BOQ = 180^\circ - C$. Karena $O_1D \perp OB$ dan $O_3D \perp OQ$, diperoleh $\angle O_1DO_3 = 180^\circ - \angle BOQ = C$.

    Di sisi lain, $\angle AO_1O = 2\angle ABO = 180^\circ - 2C$ menyebabkan $\angle O_1AO = \angle AOO_1 = C$ pada segitiga sama kaki $O_1AO$, sehingga $\angle O_1OO_3 = 180^\circ - C$.

    Jadi $\angle O_1DO_3 + \angle O_1OO_3 = 180^\circ$, artinya $OO_1O_3D$ siklis, dan $\angle O_1AO_3 = \angle O_1AO = C = \angle O_1DO_3$. Klaim 1 terbukti.
\end{proof}

\newpage
\begin{proof}[\textbf{Klaim 2.}] $AO_3 = O_3D$ jika dan hanya jika $A$, $O_1$, $D$ segaris.\\
    \textbf{Bukti Klaim.} ($\Leftarrow$) Jika $A$, $O_1$, $D$ segaris dengan $O_1$ di antara $A$ dan $D$, sudut $\angle O_3AD$ dan $\angle O_3DA$ berimpit dengan $\angle O_1AO_3$ dan $\angle O_1DO_3$, sehingga $\angle O_3AD = C = \angle O_3DA$ dan $\triangle O_3AD$ sama kaki dengan $AO_3 = O_3D$.

    ($\Rightarrow$) Jika $AO_3 = O_3D$, Aturan Sinus pada $\triangle O_1AO_3$ dan $\triangle O_1DO_3$ menyebabkan $\sin\angle AO_1O_3 = \sin\angle DO_1O_3$.

    Karena $O_1O_3$ garis sumbu $OP$, diperoleh $\angle OO_1O_3 = \angle OBP = 90^\circ - A$, sehingga
    \[
        \angle AO_1O_3 = \angle AO_1O + \angle OO_1O_3 = (180^\circ - 2C) + (90^\circ - A) = 270^\circ - A - 2C;
    \]
    karena $A > C$, diperoleh
    \[
        \angle DO_1O_3 = \angle DO_1O - \angle OO_1O_3 = (90^\circ - C) - (90^\circ - A) = A - C.
    \]

    Persamaan sinus tersebut menyebabkan $\angle AO_1O_3 = \angle DO_1O_3$ (yakni $2A + C = 270^\circ$, kontradiksi untuk segitiga lancip) atau $\angle AO_1O_3 + \angle DO_1O_3 = 180^\circ$, yang berarti $A$, $O_1$, $D$ segaris karena $A$ dan $D$ berseberangan terhadap garis $O_1O_3$. Klaim 2 terbukti.
\end{proof}
\hrulefill

Dari Klaim 2, terbukti bahwa $AO_3 = O_3D$ jika dan hanya jika titik $A$, $O_1$, $D$ segaris. \qed

\begin{remark*}
    Dari perhitungan di Klaim 2, $\angle AO_1O_3 + \angle DO_1O_3 = 270^\circ - 3C$. Jadi kedua kondisi di soal sebenarnya sama-sama setara dengan $\angle ACB = 30^\circ$. Gambar di atas memakai $A = 76^\circ$, $B = 74^\circ$, $C = 30^\circ$.
\end{remark*}

\newpage
\section{Marking Scheme yang diajukan (Bisa berbeda dengan yang dipakai juri, ya saya bukan jurinya)}
Total poin yang saya asumsikan adalah 7 poin bilangan bulat non negatif, bisa jadi berbeda dengan yang dipakai juri.
\begin{enumerate}[1)]
    \item Mengidentifikasi garis-garis sumbu dan kolinearitas dasar: membuktikan $A$, $O$, $O_3$ segaris (Lemma), $O_1D$ garis sumbu $OB$, serta $O_2O_3$ (yakni $O_3D$) garis sumbu $OQ$. (+2 poin)
    \item Membuktikan segiempat $OO_1O_3D$ siklis dan kesamaan sudut alas. (+2 poin)
    \begin{enumerate}[(a)]
        \item Menghitung $\angle BOQ = 180^\circ - C$ sehingga $\angle O_1DO_3 = C$, serta $\angle O_1AO = \angle AOO_1 = C$. (+1 poin)
        \item Menyimpulkan $OO_1O_3D$ siklis dan memperoleh hubungan kunci $\angle O_1AO_3 = \angle O_1DO_3 = C$. (+1 poin)
    \end{enumerate}
    \item Membuktikan arah ($\Leftarrow$): dari $A$, $O_1$, $D$ segaris ke $AO_3 = O_3D$ via $\triangle O_3AD$ sama kaki. (+1 poin)
    \item Membuktikan arah ($\Rightarrow$) dan eliminasi kasus degenerasi. (+2 poin)
    \begin{enumerate}[(a)]
        \item Menurunkan $\sin\angle AO_1O_3 = \sin\angle DO_1O_3$ dari $AO_3 = O_3D$ via Aturan Sinus. (+1 poin)
        \item Menghitung $\angle AO_1O_3 = 270^\circ - A - 2C$ dan $\angle DO_1O_3 = A - C$ (memakai $A > C$), lalu mengeliminasi kasus $2A + C = 270^\circ$ memakai syarat segitiga lancip sehingga $A$, $O_1$, $D$ segaris. (+1 poin)
    \end{enumerate}
\end{enumerate}

\textit{Catatan: hanya mengutip kolinearitas $A$, $O$, $O_3$ tanpa bukti maksimal 1 poin untuk bagian 1. Mengklaim $\triangle O_1AO_3 \cong \triangle O_1DO_3$ langsung dari SSA tanpa analisis kasus suplemen: $-2$ poin. Lupa menjelaskan mengapa $2A + C = 270^\circ$ kontradiksi: $-1$ poin.}
